% P19-E2: two primary experimental records, with a separate logical benchmark.
\chapter{Conditional Reasoning and the Suppression of Inferences}
\label{sec:p19-conditional-reasoning-suppression}

People sometimes endorse a conditional inference less often when another
premise is supplied. The empirical question is what changed in their
interpretation or response process. It is different from the mathematical
question whether adding a premise can invalidate classical modus ponens.

The two-study comparison here is fixed to Byrne's Experiment~1
\cite[pp.~66--69, table~1 and appendix p.~80]{Byrne1989Suppression}
and Cariani and Rips's Experiment~1
\cite[Secs.~2.3, 5.1--5.2]{CarianiRips2017Suppression}.
The latter appeared online in 2016 and in the 2017 journal volume.
The other experiments in those papers, their summaries of earlier studies,
and the wider conditional-reasoning literature are not pooled into this
comparison. These studies were selected for their common suppression
question and their differing tasks and interpretation accounts. They do
not constitute a systematic review or a population prevalence estimate.

\section{The task and its logical benchmark}
\label{sec:p19-suppression-benchmark}

\begin{hbdefinition}{def:p19-conditional-task}{An inference response and suppression}
A conditional-reasoning task supplies conditional and categorical premises,
then requests either a conclusion or an assessment of a proposed conclusion.
An \emph{endorsement rate} is the proportion of eligible responses accepting
a specified conclusion under the source's coding rule. \emph{Suppression}
here means lower endorsement in a condition with extra information than
in a specified comparison condition. It is an observed response contrast;
it does not by definition locate a mental mechanism or establish that the
same interpreted propositions were held fixed.

For the classical benchmark let \(p,q,r\) be propositional atoms and
\(\to\) the material conditional of Section~\ref{sec:prop-syntax-semantics-02}.
Here \(0\) and \(1\) code false and true, respectively. The four
traditional patterns are
\[
\begin{array}{c|c|c}
 \text{Pattern}&\text{Premises}&\text{Proposed conclusion}\\\hline
 \text{modus ponens (MP)}&p\to q,\ p&q\\
 \text{modus tollens (MT)}&p\to q,\ \neg q&\neg p\\
 \text{affirming the consequent (AC)}&p\to q,\ q&p\\
 \text{denying the antecedent (DA)}&p\to q,\ \neg p&\neg q
\end{array}
\]
An English conditional is not assumed to have that interpretation merely
because its surface wording contains an if-clause.
\end{hbdefinition}

\begin{hblemma}{thm:p19-conditional-benchmark}{proved=annex:p19.conditional-benchmark}{Validity with fixed premises and meanings}
MP and MT are valid in classical propositional semantics; AC and DA are
not valid without additional premises. For sets \(\Gamma,\Delta\) of formulas,
classical consequence is monotonic:
if \(\Gamma\entails q\), then \(\Gamma\cup\Delta\entails q\).
Consequently adding \(r\to q\), \(r\to\neg q\), or
\((p\land r)\to\neg q\) to \(p\to q,p\) cannot remove its entailment
of \(q\). The disjunction \(r\lor q\) also follows. These conclusions
concern fixed valuations and fixed formula meanings.
\end{hblemma}

The premise interpretation is therefore part of the empirical benchmark.
An everyday prediction may convey a default with background requirements;
a categorical fact may concern another time; or a listener may revise the
range of situations under consideration. These possibilities change what
the person takes the premise to assert. They do not supply a countervaluation
to material modus ponens. Conversely, an interpretation-based explanation
must be supported by evidence; calling every response a different
interpretation would insulate the account from a meaningful test.

\section{Study one: additional and alternative conditions}
\label{sec:p19-byrne-suppression-study}

\paragraph{Question, sample and design.}
Byrne tested whether an extra conditional with an additional requirement
would reduce MP and MT responses, while a conditional giving an alternative
route to the consequence would reduce AC and DA responses. The reported
analysis sample comprised 24 participants from Trinity College, University
of Dublin, randomly assigned to three groups of eight. None had received
logic tuition. Four participants were replaced for failing to complete more than
three-quarters of the arguments; the source says replacement occurred before
data analysis. Recruitment dates, age and sex distributions, and payment
are not specified in this experiment's report. The numbers recruited,
replaced and analyzed must not be conflated.

Each participant received 12 items: three contents for each of the four
inference patterns. Premise condition was between participants; inference
type and content were repeated within participants. Each booklet used a
different random item order. Testing occurred in groups of five to ten,
at the participant's own pace, with instructions against skipping items or
changing earlier responses. Participants were told to assume the premises
true and select the conclusion they thought followed. The three choices
affirmed the target proposition, denied it, or left it undetermined.
The source explicitly notes that the instructions did not explain or stress
logical necessity \cite[pp.~67--69]{Byrne1989Suppression}.

\paragraph{Actual content and response coding.}
The three domains concerned attending a play, studying late at a library,
and getting wet. In the library items, the initial antecedent concerned
an essay and the consequence concerned late study. The alternative added
textbook reading as another reason for that consequence; the additional
condition concerned the library remaining open. Thus the surface pattern
was the same two conditionals, \(p\to q,r\to q\), but the intended
relation between the antecedents differed. This is a content description
of the source stimuli, not a new experiment. The appendix uses finishing
an essay, while introductory examples use writing one; the experiment's
record is tied to its appendix rather than silently treating every
illustrative wording as the administered item.

\paragraph{Reported outcomes.}
Table~1 of the held author copy reports these percentages:
\[
\begin{array}{l|rrrr}
 \text{Premise condition}&\mathrm{MP}&\mathrm{MT}&\mathrm{DA}&\mathrm{AC}\\\hline
 \text{Simple}&96&92&46&71\\
 \text{Alternative}&96&96&4&13\\
 \text{Additional}&38&33&63&54
\end{array}
\]
Each cell pools three item responses from each of eight participants.
It is not a sample of 24 independently assigned people. The prose gives
37.5 percent for additional MP and 12.5 percent for alternative AC,
rounded to 38 and 13 in the table. From the table's rounded percentages,
the additional-minus-simple MP contrast is approximately
\(38-96=-58\) percentage points. That calculation is not a new exact
participant count or a confidence interval.

The source reports an argument-type by inference-type interaction,
\(F(6,63)=13.5,p<.01\), and one-tailed planned comparisons for the
additional-versus-simple MP and MT contrasts, each
\(t(84)=4.32,p<.001\). Alternative premises reduced DA and AC in the
reported comparisons, whereas their MP and MT comparisons were not
significant. The reported additional-premise comparisons for DA and AC
were also not significant, despite their descriptive changes. These are
the source's tests and degrees of freedom; no participant-level reanalysis
or new uncertainty interval is supplied \cite[pp.~68--69]{Byrne1989Suppression}.

\paragraph{What this comparison supports.}
The pattern is selective. Merely adding another conditional did not lower
all inference rates: alternative premises left MP at the same reported
percentage while reducing the two fallacy responses. The study constrains
an explanation that treats all extra premises as interchangeable extra
load. It does not uniquely identify counterexample search, a stored rule,
conversational interpretation or another process. Three contents and eight
people per assigned condition give limited leverage for separating
participant variation from item-specific interpretation.

The held source supplies aggregate percentages and inferential summaries,
not participant response vectors, item-by-person counts or covariances.
Those missing records prevent an independently checked clustered interval
or a reconstructed analysis from the displayed table alone. Probability
recruitment from a wider adult population is not documented.

There is also a source-version boundary: a published errata notice is
indexed at \emph{Cognition} 34(3), p.~299 (1990), but its correction text
was not available in the acquired material. The figures above are explicitly
the held author's-copy figures; incorporation or irrelevance of that errata
has not been verified. This limitation remains in the evidence record,
rather than being silently converted into corrected data.

\section{Study two: explicit deduction instructions and changed argument forms}
\label{sec:p19-cariani-rips-suppression-study}

Cariani and Rips's Experiment~1 tested suppression with two wordings of
the conditional and four argument forms. Its analyzed sample was 31
introductory psychology students at Northwestern University participating
for course credit: 15 received present/future do--will wording and 16
past-tense indicative wording. The paper reports three exclusions for
prior college logic study across its two experiments together; it does
not allocate that exclusion count to Experiment~1. Its methods describe
the two wording groups without specifying a random-assignment mechanism.
Their item orders, in contrast, were explicitly randomized separately
for each participant \cite[Sec.~2.3]{CarianiRips2017Suppression}.

The procedure presented 48 arguments to each participant: twelve
contents in each of four forms, shown individually on a computer.
The response was a binary deductive-validity judgment. The instructions
explained that true premises must guarantee a true conclusion, rather
than requesting a confidence rating. Recruitment dates, a separate
sample-size justification, detailed participant demographics and
participant-level records are not supplied in the selected methods
and results account.

For the source's rain-and-umbrella content, let \(p\) denote rain,
\(q\) the person becoming wet, and \(r\) remembering an umbrella.
The four forms, with wording paraphrased by these symbols, were:
\[
\begin{array}{l|l|l}
 \text{Form}&\text{Premises}&\text{Conclusion}\\\hline
 \text{Basic}&p\to q,\ p&q\\
 \text{Contravening}&p\to q,\ r\to\neg q,\ p&q\\
 \text{Conjunctive}&p\to q,\ (p\land r)\to\neg q,\ p&q\\
 \text{Disjunctive}&p\to q,\ r\to\neg q,\ p&r\lor q
\end{array}
\]
All four are valid under the fixed material-conditional benchmark. The
experiment investigates endorsement of their ordinary-language versions.
In particular its contravening premise has a negated consequent, unlike
Byrne's additional premise; the two conditions are not identical stimuli.

\paragraph{Rates, tests and comparisons.}
Collapsing over the two wording groups, the source reports endorsement
on 99 percent of Basic trials, 35 percent of Contravening trials,
20 percent of Conjunctive trials, and 91 percent of Disjunctive trials.
The calculated Contravening-minus-Basic difference is approximately
\(35-99=-64\) percentage points; the Conjunctive-minus-Basic difference
is approximately \(-79\) points. Twelve trials per form per participant
were specified by the design. The rounded percentages are not converted
here into purported exact endorsement counts.

The reported argument-form effect was
\(F(3,87)=124.04,\mathrm{MS}_e=0.039,p<.001\).
The reported Tukey HSD comparisons at \(\alpha=.05\) distinguished
Basic and Disjunctive from Contravening, and Contravening from Conjunctive;
Basic versus Disjunctive was not significant. Wording yielded
\(F(1,29)=0.13,p>.10\), and its interaction with form
\(F(3,87)=0.16,p>.10\). The authors consequently combined the wording
groups for these rate summaries. Failure to detect those wording effects
does not establish equivalence of the wordings. The selected primary
text supplies no exact confidence-interval endpoints for these rate
contrasts \cite[Secs.~5.1--5.2]{CarianiRips2017Suppression}.

The comparison weakens an explanation attributing all suppression to
instructions that merely request informal confidence: suppression occurred
under explicitly deductive instructions in this sample. It does not prove
that instructions never matter, or that every participant followed the
same interpretation. The high Disjunctive rate also makes a simple account
in terms of unwillingness to accept any conclusion after a longer premise
set less adequate. Its conclusion is weaker than the Contravening
conclusion, however, so that comparison changes the response target as
well as its rate.

\section{A fully worked interpretation comparison}
\label{sec:p19-suppression-reconstruction}

\begin{hbdefinition}{def:p19-suppression-reconstruction}{Two explicit reconstructions, not fitted process models}
The following finite reconstructions isolate a logical distinction needed
to interpret the experiments. They are not claims about unobserved
participant-by-participant representations.
A state in the reconstruction is a nonempty set of classical valuations;
it accepts a formula exactly when every valuation in that set satisfies it.

For the library case let \(e\) be an essay requirement, \(o\) the library
being open, and \(s\) late study. Compare the literal premise set
\(\Gamma=\{e\to s,o\to s,e\}\) with the context-sensitive
reconstruction \(\Gamma'=\{(e\land o)\to s,e\}\).
The latter treats the second sentence as drawing attention to a missing
requirement and revises what the first conditional conveys.

For the later rain case, represent a situation by its values of
\((p,q,r)\). An initial state containing
\((0,0,0),(1,1,0)\) accepts \(p\to q\); after the factual premise
\(p\), only \((1,1,0)\) remains. A context expansion that also admits
\((1,0,1)\) yields the final state
\(J=\{(1,1,0),(1,0,1)\}\). The revised rule
\((p\land\neg r)\to q\), the disabling rule \(r\to\neg q\), and
the factual \(p\) all hold there. The original unrestricted material
\(p\to q\) does not.
\end{hbdefinition}

\begin{hblemma}{thm:p19-suppression-reconstruction}{proved=annex:p19.suppression-reconstruction}{Which premise or context changed}
The literal \(\Gamma\) entails \(s\). The reconstructed \(\Gamma'\)
entails neither \(s\) nor \(\neg s\); adding \(o\) makes it entail
\(s\). The conjunction of \(p\to q\) and \(r\to q\) is equivalent
to \((p\lor r)\to q\), and is not equivalent to
\((p\land r)\to q\). Thus the conjunctive reconstruction is a change
of interpreted content, not a logical rewriting of the original premises.
In the rain state \(J\), \(q\) is not accepted but \(r\lor q\) is
accepted. A fixed nonempty set of valuations that accepts both a material
conditional and its antecedent must accept its consequent.
\end{hblemma}

Byrne's discussion proposes that background knowledge changes how the
antecedents are integrated, and relates suppression to the availability
of alternatives to a conclusion \cite[pp.~76--79]{Byrne1989Suppression}.
Cariani and Rips distinguish a broad context set from the salient subset
used to evaluate a conditional. Their account allows pragmatic processing
to expand that subset and admit previously unconsidered cases. They
explicitly do not require a person to enumerate individual possible
worlds or literally draw their diagrams
\cite[Secs.~3.1--4.2, 5.1--5.2]{CarianiRips2017Suppression}.
The two-state reconstruction illustrates one relevant expansion; it is
not a replacement for all the principles of their theory.

Both descriptions make premise interpretation consequential, but they
organize the proposed psychological work differently. Response rates alone
do not reveal whether a particular participant revised an antecedent,
retrieved an exception, changed a contextual restriction, followed a
content-sensitive rule, or used another strategy. A theory must specify
what prompts the change and predict distinctions beyond the already
observed reduction in endorsement.

\section{What the two studies jointly support}
\label{sec:p19-suppression-comparison}

\begin{center}
\begin{tabularx}{\textwidth}{@{}l X X@{}}
\toprule
Feature & Byrne, Experiment 1 & Cariani--Rips, Experiment 1\\\midrule
Analyzed people & 24; three assigned groups of eight & 31; wording groups of 15 and 16\\
Within-person responses & 12 items; three contents per inference & 48 items; twelve contents per form\\
Response task & three possible conclusions; truth assumption & binary validity judgment; necessity explained\\
Critical extra information & additional requirement with the same consequent & contravening condition with a negated consequent\\
Selected MP contrast & about 58 percentage points lower & about 64 points lower for Contravening\\
Preserved limiting comparison & alternative premises did not suppress MP & Disjunctive close to Basic; wording effects not detected\\
\bottomrule
\end{tabularx}
\end{center}

These contrasts concern different stimuli, different response formats,
different repeated-item structures and different university samples. Their
magnitudes should not be averaged as though they estimate one common
population effect. The cross-study comparison does not isolate an
instruction effect because instructions were not the only changed feature.
Within the later experiment, the nonsignificant wording tests also lack
the precision statement needed for a claim of practical equivalence.

The defensible conclusion is narrower and useful: in these tasks and
samples, adding semantically relevant information changed endorsement,
and the pattern depended on the kind of addition and the conclusion
requested. The later study shows that explicit deduction instructions did
not eliminate the pattern. Neither study refutes truth-preserving MP,
establishes one exclusive cognitive mechanism, or estimates performance
across cultures, ages, expertise levels or ordinary decisions generally.
Neither reports new participants' actual mental states for the handbook
reconstructions above.

The two-study corpus, source versions, stated exclusions and missing records
remain part of this conclusion. Raw response matrices, coding of omissions,
and a resolved original errata text would strengthen reproducibility.
Tests separating interpretation accounts would additionally need predictions
about changed instructions, elicited interpretations or counterexample
availability while preserving the other task features. Those are requirements
for further evidence, not completed experiments reported here.
